\section{Conclusions}

Following the Hill sphere transit of Beta Pic b, we have carried out an analysis on the data used in \citet{Kenworthy_2021} to search for any material transiting the star in the Hill sphere transits of Beta Pic c.

ASTEP, bRing and BRITE covered \pic c during two primary transits at MJD 58210 and 59413 and a secondary transit at 58707.
%
All transits and additional discontinuous coverages spaning from 03-02-2017 to 17-01-2023 were fitted to a model of the 1981 event to search for compelling evidence of a similar detection, however none was found.
%
Fluctuations with a magnitude increase to $a \sim$ 0.01 over a span 5-15 day by ASTEP and bRing were detected, but a lack of similar detections in the BRITE data suggest terrestrial sources of photometric noise that was not removed by the photometric pipelines.

Moreover, the circumplanetary disk's tilt and inclination was probed at 0.3 and 0.6$r_{\text{Hill}}$ using a model from \cite{Kenworthy_2021}.
%
The model's sensitivity to the midpoint of the offset was analysed, showing partial detections to the correct orientation for up to three days offset.
%
Taking the sensitivity of the model into account, the optical depth at the primary transits MJD 58210 and MJD 59413 at 0.3 and 0.6$r_{\text{Hill}}$ were computed.
%
The former transit has a coverage from all instruments and a disk bias towards inclinations smaller than 30\textdegree \ and inclinations larger than 60\textdegree.
%
bRing has the highest SNR but with the largest systematics, while ASTEP and BRITE have low SNR and positive optical depth.
%
The latter transit was only covered by bRing with a bias towards larger inclinations.
%
The results are inconsistent for the different instruments at MJD 58219 and inconsistent among both transits.

The poorly constrained midpoints of the transits and the highly sensitive model yield inconclusive results. 
%
More astrometric measurements of Beta Pic c will constrain the orbital parameters and enable an analysis of the first Hill sphere transit, along with a more general algorithm to explore the parameter space.

With an orbital period of 3.2 years, additional transits of the Hill sphere can be monitored and circumplanetary material searched for.
%
Most notably, the PLATO satellite \citep{2025ExA....59...26R}, due for launch at the end of 2026, will monitor Beta Pictoris as part of its LOPS2 field for at least two years. 
%
The next transit of the Beta Pic c Hill sphere will therefore be seen in two broadband colours with the PLATO array with space satellite millimagnitude precision, enabling the detection of material at much lower optical depths around this planet.